\documentclass[12pt]{article}
\usepackage[utf8]{inputenc}
\usepackage[T1]{fontenc}
\usepackage{hyperref}
\title{Swahili Language Resources for K-12 Curriculum Development}
\author{Compiled by Claude Sonnet 5.5 (claude-sonnet-5-5) for Andr\'as Kornai}
\date{}
\begin{document}
\maketitle

\section{Introduction}
Swahili (Kiswahili, ISO 639-3 swh) is the national language of Tanzania and is used in education there. The searches done here found little directly usable curriculum material. They did not surface a current list of Taasisi ya Elimu Tanzania (TET) primary textbooks, and I did not verify any official online dictionary portal. The resources below are the ones that appeared in search results.

\section{Primary Educational Resources}
\begin{itemize}
\item Columbia University Libraries Swahili page: lists Swahili resources, including the African Storybook (Saide) of children's stories. \url{https://library.columbia.edu/libraries/global/virtual-libraries/african_studies/languages/swahili.html}
\item University of Dar es Salaam study of Kiswahili textbooks for children in Nyamagana District (2010): reports a shortage of quality Kiswahili textbooks. \url{https://libraryrepository.udsm.ac.tz/items/92d865fb-a71c-49bc-9c79-ae0218c2ebf1}
\end{itemize}

\section{Linguistic and Reference Materials}
\begin{itemize}
\item Evaluation of A-level Kiswahili textbooks published under the 1991 textbook policy (University of Dar es Salaam): useful for judging textbook quality and gaps. \url{https://libraryrepository.udsm.ac.tz/items/599be783-4311-40f4-a597-09b826ecef6e}
\item Kioo cha Lugha article on grammar content in TET Form I and Form III Kiswahili books. \url{https://journals.udsm.ac.tz/index.php/kcl/article/view/6138/4969}
\item Mulika journal article on the role of TET in curricula. \url{https://journals.udsm.ac.tz/index.php/mj/article/view/5648/4736}
\item Lexikos review of Kamusi ya Kiswahili Sanifu, 3rd edition (2013). \url{https://lexikos.journals.ac.za/pub/article/view/1273/786}
\item TUKI English-Swahili Dictionary (Kamusi ya Kiingereza-Kiswahili), library catalogue record; print only as far as the results show. \url{https://library.kab.ac.ug/Record/2891/Details}
\item SwahiLib: Kamusi ya Kiswahili, a third-party offline app with over 17,000 entries (not an official TUKI or BAKITA product). \url{https://apps.apple.com/app/id6446771678}
\end{itemize}

\section{Religious and Community Resources}
\begin{itemize}
\item Biblia Habari Njema (BHN), read online on YouVersion. \url{https://www.bible.com/versions/393}
\item Scripture Earth page for Swahili: lists scripture texts and audio. \url{https://www.scriptureearth.org/00eng.php?iso=swh}
\item Scripture Index for Swahili on Scripture Earth. \url{https://scriptureearth.org/00i-Scripture_Index.php?iso=swh}
\end{itemize}

\section{Digital and Audio Resources}
Scripture Earth points to audio Bible recordings and to Global Recordings Network audio in several Swahili dialects. The Helsinki Corpus of Swahili and Swahili audio-learning courses were not found in these searches.

\section{Curriculum Development Context}
Swahili is a language of instruction in Tanzania, so a curriculum would build on TET syllabi and textbooks. I could not confirm a list of current titles. The search results suggest that quality textbooks are in short supply and that TET textbooks have content problems. The TET website and the Ministry of Education approved-book lists should be consulted directly.

\end{document}